% ============================================================
%  Chapter 4 -- Complementation of Omega-Regular Languages
% ============================================================

\chapter[Complementation of omega-Regular Languages]{Complementation of $\omega$-Regular Languages}
\label{ch:omega-complementation}

\aside{For finite words, complementing a
DFA is easy.  For B\"uchi automata, this
is where the theory becomes genuinely
subtle.}

Model checking naturally asks an inclusion question:
\[
  L(A_{\mathrm{sys}})\subseteq L(A_{\mathrm{spec}}).
\]
As in the finite-word case, we would like to rewrite it as
\[
  L(A_{\mathrm{sys}})\cap \overline{L(A_{\mathrm{spec}})}=\emptyset.
\]
Intersection and emptiness are manageable for B\"uchi automata.  The
hard part is complementation.

This chapter explains the algebraic proof idea used in the course:
finite-index congruences, saturation, and a Ramsey-style decomposition
of infinite words.


% ============================================================
\section{Why the Finite Trick Fails}
\label{sec:buchi-complement-fails}
% ============================================================

For finite words, we determinize, complete, and swap final states.  For
B\"uchi automata this route fails because deterministic B\"uchi automata
are strictly weaker than non-deterministic B\"uchi automata.

\begin{example}[Eventually always $a$]
\label{ex:det-buchi-gap}
Over $\Sigma=\{a,b\}$, the language
\[
  L=\{\alpha : \text{only finitely many }b\text{ occur in }\alpha\}
\]
is easy to recognize non-deterministically: guess a point after which no
more $b$ appears, then stay in an accepting state while reading only
$a$'s.  A deterministic B\"uchi automaton cannot make such a successful
guess at the right finite time.
\end{example}

The lesson is not that complementation is impossible.  The class of
$\omega$-regular languages is closed under complement.  The lesson is
that the proof cannot be the finite-word proof.


% ============================================================
\section{Congruences and Saturation}
\label{sec:congruence-saturation}
% ============================================================

The algebraic route partitions finite words into finitely many classes,
then uses products of classes to describe infinite words.

\begin{definition}[Congruence]
\label{def:congruence}
An equivalence relation $\equiv$ on $\Sigma^*$ is a
\textbf{congruence} if
\[
  u\equiv v \text{ and } u'\equiv v'
  \quad\Longrightarrow\quad
  uu'\equiv vv'.
\]
It has \textbf{finite index} if it has finitely many equivalence classes.
\end{definition}

Congruence means that only the class matters when concatenating.  If two
blocks are interchangeable now, they remain interchangeable in any
larger finite context.

\begin{definition}[Saturation]
\label{def:saturation}
Let $\equiv$ be a finite-index congruence on $\Sigma^*$.  An
$\omega$-language $L\subseteq\Sigma^\omega$ is \textbf{saturated} by
$\equiv$ if, for every pair of equivalence classes $U,V$, the language
\[
  U V^\omega
\]
is either entirely contained in $L$ or disjoint from $L$.
\end{definition}

This is a strong finiteness property.  It says that once we know the
class of a finite prefix and the class repeated forever, the membership
of the resulting infinite word is fixed.

\begin{proposition}[Saturation and complement]
\label{prop:saturation-complement}
If a language $L$ is saturated by a finite-index congruence $\equiv$,
then its complement $\overline{L}$ is saturated by the same congruence.
\end{proposition}

\begin{proof}
Fix two equivalence classes $U,V$.  Since $L$ is saturated,
$UV^\omega$ is either contained in $L$ or disjoint from $L$.  In the
first case $UV^\omega$ is disjoint from $\overline{L}$; in the second it
is contained in $\overline{L}$.  Thus the same all-or-nothing condition
holds for the complement.
\end{proof}

\begin{lemma}[Complement from saturation]
\label{lem:complement-from-saturation}
If $L$ is saturated by a finite-index congruence $\equiv$, then
$\overline{L}$ is $\omega$-regular.
\end{lemma}

\begin{proof}
First note that every equivalence class of a finite-index congruence is
a regular finite-word language.  Indeed, build a DFA whose states are
the equivalence classes, whose initial state is $[\varepsilon]$, and
whose transition is $[u]\xrightarrow{a}[ua]$.  Congruence makes this
transition well-defined.

There are only finitely many pairs of classes $(U,V)$.  Since $L$ is
saturated, each product $UV^\omega$ is either contained in $L$ or
disjoint from $L$.  Therefore
\[
  \overline{L}
  =
  \bigcup\{UV^\omega : UV^\omega\cap L=\emptyset\}.
\]
This is a finite union.  Each $U$ and $V$ is regular, so each
$UV^\omega$ is $\omega$-regular, and finite unions of $\omega$-regular
languages are $\omega$-regular.
\end{proof}

So complementation reduces to finding a finite congruence that saturates
the language of a B\"uchi automaton.


% ============================================================
\section{The Automaton Congruence}
\label{sec:automaton-congruence}
% ============================================================

Let $A=(Q,\Sigma,\Delta,q_0,F)$ be a B\"uchi automaton.  A finite word
can be summarized by what it does to pairs of states.

\begin{definition}[Path profile notation]
\label{def:path-profile-notation}
For $p,q\in Q$ and $u\in\Sigma^*$, write
\[
  p \xrightarrow{u} q
\]
when there is a path from $p$ to $q$ labelled $u$.  Write
\[
  p \xrightarrow{u}_{F} q
\]
when there is such a path that visits at least one accepting state in
$F$.
\end{definition}

The second relation is the one that finite-word automata do not need.
For B\"uchi acceptance, a repeated block is useful only if it can be
traversed while seeing acceptance.

\begin{definition}[Transition profile]
\label{def:transition-profile}
For a finite word $u\in\Sigma^*$ and states $p,q\in Q$, record two
facts:
\begin{itemize}[nosep]
  \item whether there exists a path from $p$ to $q$ labelled $u$;
  \item whether there exists such a path that visits some state in $F$.
\end{itemize}
Two words are equivalent, written $u\approx_A v$, if they have the same
answers for every pair $(p,q)$.
\end{definition}

The relation $\approx_A$ has finite index because each word determines a
finite Boolean table over $Q\times Q$.  It is also a congruence: the
transition profile of $uv$ is obtained by composing the profiles of $u$
and $v$.

\begin{proposition}[Index bound]
\label{prop:double-tilde-index}
If $A$ has $n$ states, then $\approx_A$ has at most $3^{n^2}$
equivalence classes.
\end{proposition}

\begin{proof}
Fix a finite word $u$ and a pair of states $(p,q)$.  With respect to
that pair, there are three relevant statuses:
\begin{enumerate}[nosep]
  \item no path labelled $u$ goes from $p$ to $q$;
  \item some path labelled $u$ goes from $p$ to $q$, but no such path
        visits an accepting state;
  \item some path labelled $u$ goes from $p$ to $q$ and visits an
        accepting state.
\end{enumerate}
There are $n^2$ ordered pairs $(p,q)$.  A word therefore determines a
matrix with $n^2$ entries, each entry chosen from at most three statuses.
There are at most $3^{n^2}$ such matrices.  Two words with the same
matrix are $\approx_A$-equivalent, so the number of equivalence classes
is at most $3^{n^2}$.
\end{proof}

\begin{example}[What the profile remembers]
\label{ex:profile-remembers}
Suppose $Q=\{p,q,r\}$ and $F=\{r\}$.  For a block $u$, the profile does
not remember the exact path labelled by $u$.  It remembers whether $u$
can take $p$ to $q$, $p$ to $r$, and so on, and whether an accepting
state was seen along the way.
\end{example}

This is exactly the information needed to reason about repeating finite
blocks inside an infinite run.

\begin{formaldetails}[Replacement principle]
If $u\approx_A v$, then every finite path fact about $u$ that matters
for B\"uchi acceptance also holds for $v$ with the same endpoints:
\[
  p\xrightarrow{u}q \Longleftrightarrow p\xrightarrow{v}q,
  \qquad
  p\xrightarrow{u}_{F}q \Longleftrightarrow p\xrightarrow{v}_{F}q.
\]
Thus a run built from $u$-blocks can be imitated block by block with
$v$-blocks, preserving the accepting visits that make the infinite run
B\"uchi-accepting.  This is the local reason why transition profiles
saturate $\omega$-regular languages.
\end{formaldetails}


% ============================================================
\section{Ramsey Decomposition}
\label{sec:ramsey-decomposition}
% ============================================================

The remaining difficulty is that an arbitrary infinite word need not be
literally of the form $uv^\omega$.  Ramsey's theorem lets us cut it into
blocks that are equivalent to each other.

\begin{figure}[ht]
\centering
\begin{tikzpicture}[
  cut/.style={circle, fill=teal!65!black, inner sep=1.8pt},
  arr/.style={-Stealth, thick},
]
  \draw[thick] (0,0) -- (8.6,0);
  \foreach \x/\lab in {0.5/$i_0$,2.4/$i_1$,4.3/$i_2$,6.2/$i_3$,8.1/$i_4$} {
    \node[cut] at (\x,0) {};
    \node[below=3mm] at (\x,0) {\lab};
  }
  \draw[arr, orange!80!black] (0.5,0.35) to[bend left=22] node[above] {$V$} (2.4,0.35);
  \draw[arr, orange!80!black] (2.4,0.35) to[bend left=22] node[above] {$V$} (4.3,0.35);
  \draw[arr, violet!75!black] (0.5,0.95) to[bend left=22] node[above] {$V$} (4.3,0.95);
  \node[font=\small, align=center] at (4.4,-0.9)
    {Homogeneity makes both adjacent blocks and their merge have the same class.};
\end{tikzpicture}
\caption{Ramsey homogeneity: all selected segments carry the same congruence class.}
\label{fig:ramsey-homogeneous-cuts}
\end{figure}

\begin{lemma}[Homogeneous decomposition]
\label{lem:homogeneous-decomposition}
Let $\equiv$ be a finite-index congruence on $\Sigma^*$.  Every
$\omega$-word $\alpha$ can be decomposed as
\[
  \alpha = u_0u_1u_2\cdots
\]
where $u_0\in\Sigma^*$, each $u_i$ for $i\geq 1$ is non-empty, and all
blocks $u_i$ with $i\geq 1$ belong to the same $\equiv$-class $V$.
Moreover, the class can be chosen idempotent:
\[
  V\cdot V \subseteq V.
\]
\end{lemma}

\begin{proof}
Colour each pair of positions $(i,j)$ with the $\equiv$-class of the
finite segment $\alpha[i,j]$.  Since there are finitely many classes,
the infinite Ramsey theorem gives an infinite set of positions
\[
  i_0<i_1<i_2<\cdots
\]
such that all segments $\alpha[i_m,i_n]$ have the same colour.  Taking
$u_0=\alpha[0,i_0]$ and $u_k=\alpha[i_{k-1},i_k]$ gives the desired
decomposition.

To see idempotence, take any two words $x,y\in V$.  Since $V$ is an
equivalence class and $\equiv$ is a congruence, the class of $xy$ is the
same as the class of a concatenation of two consecutive homogeneous
segments:
\[
  \alpha[i_0,i_1]\alpha[i_1,i_2]=\alpha[i_0,i_2].
\]
But $\alpha[i_0,i_2]$ has the same homogeneous colour $V$.  Hence
$xy\in V$, so $V\cdot V\subseteq V$.
\end{proof}

The point is conceptual: although the word may not be periodic, it can
be partitioned into blocks that are indistinguishable by the finite
congruence, and the repeated block class is closed under concatenating
its own blocks.


% ============================================================
\section{Closure under Complement}
\label{sec:omega-complement-theorem}
% ============================================================

\begin{theorem}[Complementation of B\"uchi-recognizable languages]
\label{thm:buchi-complementation}
If $L\subseteq\Sigma^\omega$ is $\omega$-regular, then
$\overline{L}$ is $\omega$-regular.
\end{theorem}

\begin{proof}
Let $A$ be a B\"uchi automaton for $L$, and use the finite-index
congruence $\approx_A$ from \cref{def:transition-profile}.  The Ramsey
decomposition lemma lets every infinite word be represented by a finite
prefix followed by infinitely many non-empty blocks from one
$\approx_A$-class.

It remains to justify saturation.  Consider a class product $UV^\omega$
that intersects $L$.  Choose an accepted word in it and cut it into
$u_0u_1u_2\cdots$ with $u_0\in U$ and $u_i\in V$ for $i\geq 1$.  In an
accepting run, write
\[
  q_0=s_0 \xrightarrow{u_0} s_1
      \xrightarrow{u_1} s_2
      \xrightarrow{u_2} \cdots
\]
for the states at the block boundaries.  Infinitely many of these block
segments visit an accepting state; otherwise the whole run would not be
B\"uchi-accepting.

Now take any other word
\[
  \beta=v_0v_1v_2\cdots\in UV^\omega,
\]
with $v_0\in U$ and $v_i\in V$ for $i\geq 1$.  Since $u_0$ and $v_0$
have the same transition profile, $v_0$ can take $q_0$ to the same
boundary state $s_1$.  Since each $u_i$ and $v_i$ have the same profile,
each segment $v_i$ can take $s_i$ to $s_{i+1}$, and whenever the
original $u_i$-segment visited $F$, the replacement $v_i$-segment can be
chosen to visit $F$ as well.  Stitching these replacement paths together
gives an accepting run of $A$ on $\beta$.

Therefore every word in $UV^\omega$ is accepted.  We have shown that
each class product either misses $L$ or is contained in $L$, so $L$ is
saturated by $\approx_A$.

By \cref{lem:complement-from-saturation}, the complement is a finite
union of class products $UV^\omega$, hence it is $\omega$-regular.
\end{proof}

\begin{remark}
This proof establishes effectiveness in principle: the congruence has
finitely many classes and the complement can be represented as a finite
union of $\omega$-regular components.  Practical complementation
constructions are intricate and can be expensive, but the closure result
is enough for decidability of inclusion.
\end{remark}

\begin{formaldetails}[Effective complement construction]
The proof suggests a finite procedure.  Enumerate the finitely many
transition-profile classes of $\approx_A$.  For each pair of classes
$(U,V)$, decide whether
\[
  UV^\omega\cap L(A)=\emptyset.
\]
This is a B\"uchi emptiness question, because $U$ and $V$ are regular
finite-word languages and $UV^\omega$ is $\omega$-regular.  The
complement is the finite union of exactly those products whose
intersection with $L(A)$ is empty.
\end{formaldetails}

\begin{proposition}[Ultimately periodic fingerprint]
\label{prop:ultimately-periodic-fingerprint}
Let $L_1,L_2\subseteq\Sigma^\omega$ be $\omega$-regular languages.  If
$L_1$ and $L_2$ contain the same ultimately periodic words, then
$L_1=L_2$.
\end{proposition}

\begin{proof}
Suppose not.  Then either $L_1\setminus L_2$ or $L_2\setminus L_1$ is
non-empty.  By closure under complement and intersection, that
difference is $\omega$-regular.  Every non-empty $\omega$-regular
language contains an ultimately periodic word, so the difference
contains an ultimately periodic word.  This contradicts the assumption
that $L_1$ and $L_2$ agree on all ultimately periodic words.
\end{proof}


% ============================================================
\section{Model Checking Consequence}
\label{sec:buchi-model-checking-consequence}
% ============================================================

The main verification payoff is the decidability of language inclusion.

\begin{corollary}[B\"uchi inclusion]
\label{cor:buchi-inclusion}
Given B\"uchi automata $A$ and $B$, it is decidable whether
\[
  L(A)\subseteq L(B).
\]
\end{corollary}

\begin{proof}
Construct an automaton for $\overline{L(B)}$, then an automaton for
$L(A)\cap\overline{L(B)}$, and finally test B\"uchi emptiness by looking
for a reachable accepting cycle.
\end{proof}

This is the automata-theoretic backbone of model checking: search for an
execution of the system that violates the specification.  If none
exists, inclusion holds.

\begin{summarybox}
\begin{itemize}[nosep]
  \item B\"uchi complementation cannot be done by simply swapping final states.
  \item Finite-index congruences make infinite behaviours finitely classifiable.
  \item Saturation turns complementation into a finite union construction.
  \item Ramsey decomposition explains why arbitrary infinite words can be handled by repeated class products.
  \item Inclusion of B\"uchi languages is decidable via complement, intersection, and emptiness.
\end{itemize}
\end{summarybox}

\section*{Exercises}
\addcontentsline{toc}{section}{Exercises}

\begin{chapex}[Finite-index class automaton]
Let $\equiv$ be equivalence modulo parity of word length.  Build the DFA
whose states are the two equivalence classes and whose transitions append
one symbol.  Which class contains $\varepsilon$?
\end{chapex}

\begin{chapex}[Why profiles need accepting information]
Give two finite paths from a state $p$ back to $p$, one visiting an
accepting state and one not.  Explain why ordinary reachability
information is insufficient for B\"uchi complementation.
\end{chapex}

\begin{chapex}[Inclusion reduction]
Write the three automata operations needed to decide
$L(A)\subseteq L(B)$ for B\"uchi automata.  Which operation is the
technically hardest one?
\end{chapex}
